%
%  THESISBOEK
%
%  Dit bestand zorgt voor algemene (layout)definities, en groepeert de
%  afzonderlijke LaTeX-files tot een geheel.
%
%  @author Erwin Six, David De Reu, Brecht Vermeulen
%

\documentclass[11pt,a4paper,oneside,notitlepage]{book}
\usepackage[english]{babel}
\usepackage{a4wide}                     % Iets meer tekst op een bladzijde
\usepackage{pdfpages}					% Pdf includeren
\usepackage{subcaption}					% 2 figuren naast mekaar, met elk een caption
\usepackage[bottom]{footmisc}			% footnotes forceren naar bottom
\usepackage{graphicx}                   % Om figuren te kunnen verwerken
\usepackage[latin1]{inputenc}           % Om niet ascii karakters rechtstreeks te kunnen typen
\usepackage{fancyhdr}
\usepackage{hyperref}
\usepackage{float}

\usepackage{nomencl}
\makenomenclature

\usepackage[nottoc]{tocbibind}
\usepackage[numbers]{natbib}

% marges aanpassen
% (opmerking: moet *voor* inclusie van fancyhdr package komen)
\setlength{\hoffset}{-1in}
\setlength{\voffset}{-1in}
\setlength{\topmargin}{2cm}
\setlength{\headheight}{0.5cm}
\setlength{\headsep}{1cm}
\setlength{\oddsidemargin}{3.5cm}
\setlength{\evensidemargin}{3.5cm}
\setlength{\textwidth}{16cm}
\setlength{\textheight}{23.3cm}
\setlength{\footskip}{1.5cm}

\pagestyle{fancy}

\renewcommand{\chaptermark}[1]{\markright{\MakeUppercase{#1}}}
\renewcommand{\sectionmark}[1]{\markright{\thesection~#1}}

\newcommand{\headerfmt}[1]{\textsl{\textsf{#1}}}
\newcommand{\headerfmtpage}[1]{\textsf{#1}}

\fancyhf{}
\fancyhead[LE,RO]{\headerfmtpage{\thepage}}
\fancyhead[LO]{\headerfmt{\rightmark}}
\fancyhead[RE]{\headerfmt{\leftmark}}
\renewcommand{\headrulewidth}{0.5pt}
\renewcommand{\footrulewidth}{0pt}

\fancypagestyle{plain}{ % eerste bladzijde van een hoofdstuk
  \fancyhf{}
  \fancyhead[LE,RO]{\headerfmtpage{\thepage}}
  \fancyhead[LO]{\headerfmt{\rightmark}}
  \fancyhead[RE]{\headerfmt{\leftmark}}
  \renewcommand{\headrulewidth}{0.5pt}
  \renewcommand{\footrulewidth}{0pt}
}

\setlength{\parindent}{0cm}             % Inspringen van eerste lijn van paragrafen is niet gewenst

% Vandaar dat we expliciet aangeven wanneer we wensen dat een nieuwe paragraaf begint:
% \par zorgt ervoor dat er een nieuwe paragraaf begint en
% \vspace zorgt voor vertikale ruimte.
\newcommand{\npar}{\par \vspace{2.3ex plus 0.3ex minus 0.3ex}}

\makeatletter
\renewcommand{\@chapapp}{}% Not necessary...
\newenvironment{chapquote}[2][2em]
  {\setlength{\@tempdima}{#1}%
   \def\chapquote@author{#2}%
   \parshape 1 \@tempdima \dimexpr\textwidth-2\@tempdima\relax%
   \itshape}
  {\par\normalfont\hfill--\ \chapquote@author\hspace*{\@tempdima}\par\bigskip}
\makeatother

% Nieuw commando om figuren in te voegen. Gebruik:
\newcommand{\mijnfiguur}[4][H]{            % Het eerste argument is standaar `ht'. op H zetten voor HIER EN NERGENS ANDERS
    \begin{figure}[#1]                      % Beginnen van de figure omgeving
        \begin{center}                      % Beginnen van de center omgeving
            \includegraphics[#2]{#3}        % Het eigenlijk invoegen van de figuur (2: opties, 3: bestandsnaam)
            \caption{#4\label{#3}}          % Het bijschrift (argument 4) en het label (argument 3)
        \end{center}
    \end{figure}
    }

\graphicspath{{fig/}}               % De plaars waar latex zijn figuren gaat halen.


% PDF specifieke opties, niet strict noodzakelijk voor een thesis.
% Is hetgeen verschijnt wanneer je in acroread de documentproperties bekijkt.
\hypersetup{
    pdfauthor = {Andreas De Lille},
    pdftitle = {title},
    pdfsubject = {research to emotion recognition using EEG data},
    pdfkeywords = {EEG, Emotion recognition,BCI}
}

\newenvironment{samenvatting}{\small\itshape}{}

\begin{document}

% titelblad (voor kaft)
%  Titelblad

% Opmerking: gaat uit van een \baselinestretch waarde van 1.5 (die moet
% ingesteld worden voor het begin van de document environment)

\begin{titlepage}

\setlength{\hoffset}{-1in}
\setlength{\voffset}{-1in}
\setlength{\topmargin}{1.5cm}
\setlength{\headheight}{0.5cm}
\setlength{\headsep}{1cm}
\setlength{\oddsidemargin}{3cm}
\setlength{\evensidemargin}{3cm}
\setlength{\footskip}{1.5cm}
\enlargethispage{1cm}
% \textwidth en \textheight hier aanpassen blijkt niet te werken


\fontsize{12pt}{14pt}
\selectfont

\begin{center}

%\includegraphics[height=2cm]{fig/ruglogo}

\vspace{0.5cm}

Faculty of Engineering and Architecture\\
Departement ELIS\\
2015--2016 \\

\vspace{3.5cm}

\fontseries{bx}
\fontsize{17.28pt}{21pt}
\selectfont

Recognize Emotion in the brain using EEG Data

\fontseries{m}
\fontsize{12pt}{14pt}
\selectfont

\vspace{.6cm}

by 

\vspace{.4cm}

Andreas DE LILLE

\vspace{3.5cm}


\begin{tabular}{ll}
Promotors: & Prof.~J.~DAMBRE\\
 & Dr.~Ir.~P.~VAN MIERLO\\
Assistent: & Ir.~T.~VERHOEVEN\\
\end{tabular}
 

\vspace{3cm}


\end{center}
\end{titlepage}
 %moet nog via plato

% lege pagina (!!)

% titelblad (!!) => via pdf van plato!

% geen paginanummering tot we aan de inhoudsopgave komen
\pagestyle{empty}

% voorwoord met dankwoord en toelating tot bruikleen (ondertekend)
%%  Voorwoord (dankwoord) en toelating tot bruikleen

\newpage

\noindent \textbf{\huge Voorwoord}

\vspace{1.5cm}

\noindent
Hier komt wat tekst.

\addvspace{4cm}

\noindent David De Reu, mei 2002\newpage

\noindent \textbf{\huge Toelating tot bruikleen}

\vspace{1.5cm}

\noindent
``De auteur geeft de toelating deze scriptie voor consultatie beschikbaar
te stellen en delen van de scriptie te kopi\"eren voor persoonlijk
gebruik.\\
Elk ander gebruik valt onder de beperkingen van het auteursrecht,
in het bijzonder met betrekking tot de verplichting de bron uitdrukkelijk
te vermelden bij het aanhalen van resultaten uit deze scriptie.''

\addvspace{4cm}

\noindent David De Reu, mei 2002


% abstract
%%  Overzichtsbladzijde met samenvatting

\newpage
\addcontentsline{toc}{chapter}{Overview}
{
\setlength{\baselineskip}{14pt}
\setlength{\parindent}{0pt}
\setlength{\parskip}{8pt}

\begin{center}

\textbf{\huge
A Comparative Study of \\
Physiological Feature Selection Methods \\
for Emotion Recognition\\
}

by

Andreas DE LILLE

\end{center}

Supervisors: Prof.~J.~DAMBRE and Dr.~Ir.~P.~BUTENEERS \\
Counsellor: Ir.~T.~VERHOEVEN

Master's dissertation submitted in order to obtain the academic degree of\\
Master of Science in Computer Science Engineering

Department of Electronics and Information Systems\\
Chair: Prof. dr. ir. Rik Van de Walle\\
Faculty of Engineering and Architecture\\
Academic year 2015-2016\\



\section*{Summary}

An emerging field of research is the field of emotion recognition. Emotion can be observed in many ways, but the most reliable method is to use physiological signals.  This method uses machine learning to classify emotions based on characteristics or features extracted from the signals. To do so, good, reliable features are needed. This method compares a wide range of features and feature selection techniques to study the physiological responses triggered by emotion.

\section*{Keywords}

Emotion recognition, Physiological Signals, Machine Learning, Feature Selection

}

\newpage % strikt noodzakelijk om een header op deze blz. te vermijden


\pagestyle{fancy}
\frontmatter

% inhoudstafel
\tableofcontents

%afkortingen
\printnomenclature.

% eventueel: lijst van figuren en tabellen
%\listoffigures
%\listoftables


% opmaak voor het eigenlijke boek; onderstaande lijnen
% weglaten als de eerste regel van een nieuwe alinea moet
% inspringen in plaats van extra tussenruimte
\setlength{\parindent}{0pt}
\setlength{\parskip}{0.5\baselineskip plus 0.5ex minus 0.2ex}
\setlength{\parskip}{1ex plus 0.5ex minus 0.2ex}

% hoofdstukken
\mainmatter

% hier worden de hoofdstukken ingevoegd (\includes)
\chapter{Introduction}
{\samenvatting This chapters introduces the masterthesis. It starts by introducing the basic concepts of emotion recognition based on physiological signals, machine learning. Then it explains the problem statement and the goal of the thesis, followed by an overview of the next chapters.}

\section{Emotion recognition}

Human-to-machine communication, where humans communicate with machines or computer agents, is becoming more and more common\citep{hummaccom}. Fully understanding human communication is a complex problem. In addition to verbal communication, non-verbal communication is also used to exchange information\citep{EmotionRelativePower}. To better understand human-to-machine communication, more insight in the non-verbal communication is needed. Emotion recognition is becoming an increasingly important field as a result\citep{RealTimeEEGEmotion}.

\npar 

Emotion recognition is the proces of recognizing a subject's emotional state. In psychology a clear distinction between physiological behavior and the conscious experience of an emotion, called expression\cite{ExtendedPaper} is made. Expression consists of many parts, including facial expression, body language and voice concern\citep{EMSpeech}. Unlike expression, the physiological aspect of an emotion, e.g. heart rate, skin conductance and pupil dilation, is much harder to control. This makes emotion recognition based on physiological signals more robust to social masking\citep{PhytoEm}. Social masking is the process where an individual masks or hides their emotions to conform to social pressure. To really know one's emotions, it seems, one has to research the physiological aspect of the emotion.

\npar

\label{valarrdomspace}
Before emotions can be recognized, an objective class model describing different emotions is needed. A simple way of achieving this is using several discrete emotions, e.g. anger, joy, sad and pleasure. A more convienent model to classify emotions is the bipolar arousal-valence model\cite{ExtendedPaper,RealTimeEEGEmotion}, which places emotions in a two dimensional space. The main advantage of using a continuous multidimensional model, is that all emotions are modelled in its space, even when no particular discrete label can be used to define the current feeling. Figure \ref{ArousalValenceModel} shows the mapping of different emotions for this model. 

\npar

The valence-arousal model consists of two dimensions. Arousal indicates how active a person is and ranges from inactive/bored to active/excited. The valence indicates if the emotion is perceived as positive or negative. Even though arousal and valence describe emotions quite well, a third dimension, dominance, can also be added. This third dimension indicates how strong the emotional feeling was and ranges from a weak feeling to an empowered, overwhelming feeling. The dominance component can aid to filter out samples of strong feelings, since feelings with low dominance are less likely to show significant effects.

\mijnfiguur{width=0.45\textwidth}{ArousalValenceModel}{The arousal - valence model maps emotions in a two dimensional plane\citep{ValArrFig}}

\subsection{Physiological signals}

Now that the classification model is defined, the different physiological signals will be explained. As mentioned before, these signals are used to do automatic emotion recognition. Physiological signals can be divided in two subgroups: brain activity and other signals, e.g. heart rate, respiration rate, etc. Different technologies exist to record brain activity. The most convenient method is electroencephalography (EEG)\nomenclature{EEG}{Electroencephalography}, since it is a non-invasive method. Non-invasive methods, in contrast to invasive methods require no surgery. In case of EEG, they simply measure electrical activity using electrodes placed on the scalp.

\npar

Electrical activity in the brain is generated when an incoming signal arrives in a neuron. This triggers some sodium ions to move inside the cell, which in turn, causes a voltage rise\cite{ExtendedPaper}. When this increase in voltage reaches a threshold, an action potential is triggered in the form of a wave of electrical discharge that travels to neighbouring neurons. When this reaction occurs simultaneously in a lot of neurons, the change in electrical potential becomes significant enough, it is measured by the EEG surface electrodes. EEG can thus only capture synchronized activity of many, many neurons\cite{ExtendedPaper}. This explains why EEG signals have low spatial resolution capabilities. EEG measurements consist of electrical potentials of different channels, measured over time, like shown in Figure \ref{eegexample}.

\mijnfiguur{width=0.9\textwidth}{eegexample}{EEG measurements is a trace electrical potentials of different channels over time.\citep{EEGExampleFig}}

Signals originating from the cortex, close to the skull, are easier to measure, while signals originating deeper in the brain cannot be observed directly. Even for signals originating close to the cortex, EEG is far from precise as the bone between the the cortex and electrodes distorts the signal. Additionally, other artifacts like eye and muscle movement add a lot of noise to the signal. This explains why EEG signals are very noisy by nature. Noise removal techniques are therefor advised\citep{noiseRem}. Note that even though EEG data contains a lot of noise and has a low spatial resolution, it still provides significant insight into the electrical activity of the cortex while offering excellent temporal resolution\cite{GivenPaper}.

\npar

To ensure that experiments are replicable, standards for locations of electrodes have been developed. One of these systems is the 10/20 system, an internationally recognized method to describe the location of scalp electrodes\cite{TenTwentyManual}. The numbers 10 and 20 refer to the distances between the electrodes, which are either 10\% or 20\% of the total front-back or left-right distance of the scalp, this is depicted in Figure \ref{1020ElectrodePlacementSystem}. 

\mijnfiguur{width=0.8\textwidth}{1020ElectrodePlacementSystem}{The electrode placement of a 23 channel system\cite{1020Site}.}

Each site is identified with a letter that determines the lobe and a number that determines the hemisphere location.
\begin{itemize}
\item \textbf{F:} Frontal
\item \textbf{T:} Temporal
\item \textbf{C:} Central
\item \textbf{P:} Parietal
\item \textbf{O:} Occipital
\end{itemize}
Note that no central lobe exists; the C letter is only used for identification purposes. The letter z indicates that the electrode is placed on the central line. Even numbers are used for the right hemisphere, while odd numbers are used for the left hemisphere. Note that the 10/20 system does not require a fixed number of channels. Some experiments may use a different set of channels, but they all follow the same naming convention. In this work, a 32 channel EEG cap is used. The corresponding electrode locations are shown in Figure \ref{1020labels}

\mijnfiguur{width=0.7\textwidth}{1020labels}{Placement of the 32 electrodes in this thesis.}

In the frequency domain, brain waves are usually split up into different bands\cite{EmotionRelativePower,WavesSite}, with a different medical interpretation for each band. These wavebands \label{wavebands} are:
\begin{enumerate}
\item \textbf{Alpha:} 8-13Hz, indicate how relaxed and/or inactive the brain is.
\item \textbf{Beta:} 13-30HZ, indicate a more active and focused state of mind.
\item \textbf{Gamma:} 30-50Hz, relate to simultaneous processing of information from different brain areas.
\item \textbf{Delta:} 0-4hz, these waves are generated during dreamless sleep and meditation.
\item \textbf{Theta:} 4-8Hz, occurs during dreaming.
\end{enumerate}

\npar

Even though EEG is used in this thesis, alternative methods to measure brain activity exist. What follows is an overview of some of these techniques.
\begin{itemize}

\item Magnetoencephalography (MEG)\nomenclature{MEG}{magnetoencephalography} use magnetic fields to measure brain activity\citep{meg}. Since MEG is more prone to noise from external magnetic signals, i.e. the earth's magnetic field and electromagnetic communication, a magnetic shielded room is required, making this method very expensive and not mobile.

\item Functional magnetic resonance (fMRI) \nomenclature{fMRI}{Functional Magnetic Resonance}\citep{fMRI}: works by detecting changes in blood oxygenation and blood flow. An active area of the brain consumes more oxygen and has an increased blood flow.

\item Computed tomography (CT) \nomenclature{CT}{Computed Tomography}\citep{CT}: uses X-rays to create an image of the brain. 

\item Positron emission tomography (PET) \nomenclature{PET}{Positron Emission Tomography}\citep{PET}: this methods uses trace amounts of short-lived radioactive material. When this material undergoes decay, a positron is emitted that is picked up by a detector.

\item Near infrared spectroscopy (NIRS) \nomenclature{NIRS}{Near Infrared Spectroscopy}\citep{NIRS}: an optical technique to measure blood oxygenation in the brain. This technique works by shining light in the near infrared part of the spectrum through the skull and measuring how much remerging light is attenuated.

\end{itemize} 

In addition to brain activity, other physiological signals are also used in this work. The most known signal is the heart rate, which measures the number of contractions per minute. Respiration rate gives to number of breaths a human takes in one minute\citep{DEAP}. Another physiological signal is the galvanic skin response. The galvanic skin response measures the electrical characteristics of the skin\citep{GSR, DEAP}. In addition to the electrical characteristics, the temperature of the skin can also be measured. A plethysmograph is another physiological signal, that measures changes in volume within an organ\citep{DEAP}. A plethysmograph can be used to measure a subject's blood pressure and heart rate.

\section{Machine learning}
Machine learning is the missing link between the physiological signals and the emotion recognition. It is, in short, an input output model, that takes physiological signals and maps them to an emotional state. Machine learning is a very broad domain. As a result, this discussion will be limited to an introduction of the basic machine learning concepts with the focus on the application of machine learning and machine learning techniques used in this thesis. 

\npar

A possible definition for machine learning is: "the science of getting computers to act without being explicitly programmed"\citep{MLDef}. To do so, machine learning uses pattern recognition to find patterns or structure in the data. A simple example of machine learning is the Optical Character Recognition (OCR)\nomenclature{OCR}{Optical Character Recognition}, where a computer recognises characters in pictures\citep{OCR}. An example of OCR is shown in Figure \ref{OCR}.

\mijnfiguur{width=0.3\textwidth}{OCR}{In optical character recognition, a computer uses machine learning to find characters in an image\cite{OCRFigRef}.}

To further explain how machine learning works, have a look at the following example. Suppose one has a price list of houses that are for sale combined with their total area, shown in Table \ref{mlexampleTable}. Logic sense dictates us that a bigger house will have a higher price than a smaller house. The total area is a characteristic of the house that helps us in determining the price. In the context of machine learning, the characteristic 'total area', will be called a feature as the asking price of a house is correlated to the total area.

\begin{table}[H]
\centering
\begin{tabular}{ll}
\textbf{Area of the house ($m^{2}$)} & \textbf{Price ( x 1000 EUR)} \\
70                              & 312                          \\
73                              & 429                          \\
76                              & 174                          \\
79                              & 410                          \\
82                              & 334                          \\
$\vdots$                        & $\vdots$
\end{tabular}
\caption{total area of different houses and their corresponding asking prices.\label{mlexampleTable}}
\end{table}

One possible way of predicting the asking price of a house is machine learning. Machine learning works in several steps, first you train the machine learning algorithm with a list of asking prices and the corresponding area of the house. This process is called training or fitting and gives the machine learning component an idea to what price corresponds to a house with a certain area. Once trained, the algorithm's output will look like Figure \ref{mlexample}. The black dots represent the data points from Table\ref{mlexampleTable}. The blue line represents the predicted price for different area. The predicted price is simply defined by the total area of the house multiplied by some weight.

\mijnfiguur{width=0.9\textwidth}{mlexample}{The price of a house is determined by its total area.}

\clearpage

Even though, the blue line looks reasonable, there is sometimes a big difference between the predicted value and the actual value. This is due to the fact that the area of the house is only one feature that determines the price. Other features, like the number of bedrooms or the location of the house, were not taken into consideration. Adding additional features, gives more insight into the data, e.g. a house with 5 bedrooms is more expensive than a house with only 3 bedrooms. Having more features is thus likely to improve the performance of the machine learning algorithm.

\npar

There are many machine learning algorithms. One way to group these algorithms is to look at the produced output. In the asking price examples above, the output is a price, which is (more or less) a continuous value. Machine learning problems that require the output of a continuous value, are called regression problems\citep{prml}. In the OCR example above, a picture of a character is classified as a character. This means that OCR is a classification problem, as there are only a limited number of characters in an alphabet\citep{prml}.

\npar

Another way to group algorithms is based on their training data\citep{prml}. In the asking price examples above, the training data consists of labelled results. Labelled training data corresponds to data where the correct output (in this case the asking price) is given for each input (the area). This type of machine learning is referred to as supervised machine learning\citep{prml}. The alternative is unsupervised machine learning\citep{prml}. Unsupervised learning often results in finding groups of similar data points (clustering), without knowing the actual labels. Note that the combination of supervised and unsupervised data, known as semi-supervised learning, is also possible\citep{semiSup}. Imagine a dataset with 5000 webpages that need to be grouped into 10 distinct categories, e.g. science, nature, cooking, ... . Only 100 of the 5000 pages in the train set are labelled. An approach to solve this problem could be to first cluster the pages in similar groups using unsupervised learning. As soon as a group contains a single labelled page, all pages in the group can be labelled accordingly. This is possible because clustering returns groups of similar samples. Semi supervised learning has the advantage that one can also use unlabelled data, which is often easier and cheaper to obtain, unlike labelled data which is usually quite rare; if there was a fast and easy way to label the data then there would not be a need for machine learning.

\subsection{Over- and underfitting}

Over and underfitting is a common problem in many machine learning projects\citep{prml}. Suppose the example in Figure \ref{overunderfitting}, where one tries to find a good function to fit the given data points. Looking at the three proposed functions, one can easily see that the middle figure corresponds to the most logical generator function<\footnote{A generator function is a theoretical concept to describe the 'actual' function that generated the outputs.} of the red points.

\mijnfiguur{width=0.9\textwidth}{overunderfitting}{Overfitting versus underfitting\cite{overunderfittingFig}.}

\npar

The figure on the left corresponds to an underfit, where the proposed function is not able to capture sufficient detail of the points. The function is not complex enough to approach the generator function. As a result the best fit will always contain a relatively big error.

\npar

The function on the right corresponds to an overfit. The function fits or 'goes through' each point exactly, which will a very low error for these data points. However, one can see that the behaviour of the hypothesis function in between data points is not what one would expect. This is the result of using a too complex function to fit the data. As a results, all sample points are matched exactly, which results in a very low error for these points. However, the problems arise when the algorithm is test on unseen points. The algorithm will have a much higher error for those points. 

\npar 

Part of a good machine learning algorithm is finding the right tradeoff between overfitting and underfitting. In case of the aforementioned overfitting, it would be better to lower the performance of the algorithm on the sample datapoints, to gain performance on the 'unseen' datapoints. Different techniques exist to estimate how good an algorithm performs on unseen points. Often a part of the sample points is put in a test set that is neglected during training. After the algorithm is designed, the performance on the test set will indicate how well the algorithm generalises. It is only the performance of the test set that gives a fair estimation of the performance of a machine learning algorithm.

\subsection{Feature selection}
Feature selection is a technique that aims at selecting the features that perform well, while trying to remove irrelevant features\citep{rfPaper}. The advantages of having a smaller features set are twofold. First having fewer features, will lower the risk of overfitting\citep{rfPaper}. Second, knowing which features are important makes it possible for humans to interpret the machine learning model. In this thesis, knowing which features are relevant might help in gaining insight in the processing of emotion by the brain. 

\clearpage

There exist several approaches to do feature selection. The first one is to simple use a statistical metric and remove all features with low correlation to the output. Another approach is to look at the weights of a model. When a machine learning model gives a large weight to a feature, then that feature is considered more important than a feature with an assigned weight close to zero. Embedded methods also exist, they rely on the build in feature selection mechanisms of some machine learning algorithms. A more thorough  overview of the different feature selection techniques is given in Section \ref{FSSel}. 

\section{Problem statement}

A lot of different physiological features are reported in the literature. Unfortunately, the literature does not fully agree on a specific set of features nor does it agree on what EEG channels and/or frequency bands are most important for emotion recognition. The features that are reported in different studies are often quite different, as you can see in Table \ref{diffFeat} below.

\begin{table}[H]
\centering
\begin{tabular}{ll}
\textbf{study} & \textbf{features used}                         \\
\citep{ref4}     & Alpha and beta power                           \\
\citep{ref7}     & PSD and asymmetry features                     \\
\citep{ref8}     & PSD                                            \\
\citep{ref6}     & discrete wavelet transform of alpha, beta and gamma band \\
\citep{ExtendedPaper}	&	alpha/beta ratio, Fpz beta and alpha band power \\
\citep{killyPaper} & PSD, RCAU, DCAU, DASM, RASM, DE \\
\end{tabular}
\caption{Six different papers on emotion recognition, six different feature sets\label{diffFeat}.}
\end{table}

\npar

Another related problem with physiological signals is that they are very personal by nature. Features that work well for one person might not work well for another person\citep{DEAP}. Finding a set of features that works well for all persons is hard, but it might make the system more robust against personal differences. 

\npar

The last problem is that is hard to compare the performance of different physiological feature studies, as they do not share the same dataset.

\section{Goal of the thesis}
The first goal is finding relevant features for emotion recognition in a person specific setting. This is already quite challenging as there are fuzzy boundaries and individual variation of emotion\citep{emorecoghard}. To do so, the output of different feature selection methods is compared. In a successful scenario, good features are found. These features could be used by a machine learning algorithm to accurately predict the emotions of one person. Some attention will also be spend on comparing non-EEG and EEG features to see which whether it is useful or not to include EEG and/or non-EEG signals in the emotion recognition.

\clearpage

The second goal is finding features for emotion recognition in a cross-person setting. In this setting features should generalise well across different persons, thus the algorithm should be able to recognize emotions from unseen persons. The comparison for non-EEG and EEG features will also be done here. Emotion recognition is harder in a cross-person setting, since physiological signals are very personal\citep{DEAP}.

\npar 

Both goals are tackled by comparing a large range of different feature selection methods combined with a huge feature set. Additionally, the accuracy on the DEAP, a dataset designed to compare different emotion recognition studies\citep{DEAP} will be reported. This will ensure that the results obtained in this thesis can serve as a benchmark for future research. This is important as performance of emotion recognition algorithms based on physiological signals often varies a lot for different datasets\citep{PhytoEm}.

\npar 

The contents of this thesis are as follows. The next chapter gives an overview of the dataset, features and feature selection methods that are used in this thesis. It also gives an overview of similar state of the art emotion recognition studies.
 
\npar
Chapter 3 and 4 give an overview of the obtained results for person specific and cross-person emotion recognition respectively. Chapter 5 gives the conclusion of this work. Chapter 6 gives an overview of future research that is possible.
\chapter{Machine Learning}
{\samenvatting }


\section{What is machine learning?}


\section{train, test, validation error}

\section{svm}

\section{lda}

\section{CSP}

\section{ANOVA}

\section{Pearson correlation}

\chapter{A first look at the data}
{\samenvatting }

\section{The DEAP dataset}
This thesis uses the DEAP dataset\cite{DEAP}, a dataset for emotion analysis that is publicly available for academic research. This dataset contains EEG recordings of 32 participants, each watching 40 one minute excerpts of music videos. Each video was rated individually by each person on valence, arousal, dominance and liking. The first three ratings correspond to the valence, arousal and dominance space of an emotion. The liking component indicates how much the person liked the video excerpt and should not be confused with the valence component. The liking measure inquires about the participants' tastes and not their feelings, i.e. a person can like a video that triggers angry or sad emotions. Since liking is not a part of the emotion space, the liking rates are neglected for this thesis.

\npar

For assessment of these scales, the self-assessment manikins (SAM)\nomenclature{SAM}{self-assessment manikins}, were used\cite{DEAP}. SAM visualizes the valence, arousal and dominance scale with pictures. Each picture corresponds to a discrete value, but the user can click in between the different figures, which makes the output continuous. The used figures are shown in Figure \ref{SAM}. The first row gives the valence scale, ranging from sad to happy. The second row shows the arousal scale, ranging from bored to excited. The last row represents the different dominance levels. The left figure represents a submissive emotion, while the right figure corresponds with a dominant feeling.

\mijnfiguur{width=0.5\textwidth}{SAM}{The images used for the SAM\cite{DEAP}.}


\chapter{Methods}
{\samenvatting This chapter starts by explaining well-known features for emotion recognition. Then an overview of some emotion recognition studies is given. The chapter ends by going through frequently used feature selection methods.}

\section{Features}
\label{featuresExplained}
Usually, good features are needed to train a machine learning algorithm. In the context of this thesis, good features should be correlated with the emotion. Two categories of features are observed: EEG features and non-EEG features. Both categories are covered in the following sections.

\subsection{EEG-features}
EEG features are extracted from the electroencephalography measurements from the subject's scalp. The power spectral density (PSD) \nomenclature{PSD}{Power Spectral Density} of a signal gives the distribution of the signal's energy in the frequency domain. By calculating the spectral density for different wavebands of the signal, one can determine how much alpha, beta, ... power is in the signal.

\npar

Differential entropy (DE)\nomenclature{DE}{Differential Entropy} is defined as follows \citep{killyPaper} \\
\begin{center}
$ - \int_{-\infty}^{\infty} \frac{1}{\sqrt{2\pi\sigma^2}} exp(\frac{(x-\mu)^2}{2\sigma^2}) log(\frac{1}{\sqrt{2\pi\sigma^2}}) exp(\frac{(x-\mu)^2}{2\sigma^2})dx$
\end{center}
According to \citep{diffEnt}, the differential entropy of a certain band is equivalent to the logarithmic power spectral density for a fixed length EEG sequence, which simplifies the calculations significantly.
\begin{center}
$DE_{channel} = log(PSD_{channel})$
\end{center}

\npar

The most known feature for valence recognition is the frontal asymmetry of the alpha power\cite{GivenPaper}. The right hemisphere is generally speaking, more active during negative emotion than the left hemisphere which is in turn more active during positive emotions\cite{RealTimeEEGEmotion,EEGDatasets,killyPaper}. The asymmetry can be calculated in different ways. First, one can calculate the differential asymmetry (DASM) \nomenclature{DASM}{Differential Asymmetry}, where the left alpha power is subtracted from the right alpha power.

\begin{center}
$DASM = DE_{left} - DE_{right}$
\end{center}

Another way to measure the asymmetry is by division. The Rational Asymmetry (RASM) \nomenclature{RASM}{Rational Asymmetry} does exactly this and is given by: \\

\begin{center}
$RASM = \frac{DE_{left}}{DE_{right}}$
\end{center}

With $DE_{left}$ and $DE_{right}$ being the left and right differential entropy respectively.

\npar

Another reported feature in literature is the caudality, or the asymmetry in fronto-posterior direction\cite{caudality}, meaning the difference in power between the front and the back of the scalp. This can again be calculated in two ways. The first method is the differential Caudality (DCAU) \nomenclature{DCAU}{Differential Caudality} is defined as: \\

\begin{center}
$DCAU = DE_{front} - DE_{post}$
\end{center}

Another way to determine the Caudality is the Rational Caudality (RCAU) \nomenclature{RCAU}{Rational Caudality}, which is defined as:

\begin{center}
$RCAU = \frac{DE_{front}}{DE_{post}}$
\end{center}

With $DE_{front}$ and $DE_{post}$ being the frontal and posterior power respectively.

\npar

Arousal is usually determined, by looking at the different wavebands\citep{ExtendedPaper}. Each waveband has their own medical interpretation, see \ref{wavebands}. Alpha power corresponds to a more relaxed brain, while Beta power corresponds to a more active brain. The alpha / beta ratio therefore seems a good indicator for the arousal state of a person.

\npar

The Alpha/ Beta ratio is limited to comparing two wavebands. Other frequently used features are fractions of PSD. Where the fractions of waveband power is determined for a channel, given by:

\begin{center}
$frac_{band} = \frac{power_{band}}{power_{total}}$
\end{center}

The fractions give insight in the distributions of wavebands at different channel locations.

\subsection{non-EEG features}
The aforementioned EEG features are just one class of physiological features. The DEAP dataset contains several physiological measurements, listed below \citep{DEAP}. For each of these measurements the average, standard deviation, variation, median, minimum, maximum and the standard deviation is calculated.

\npar

The Galvanic Skin Response \nomenclature{GSR}{Galvanic Skin Response} uses two electrodes on the middle and index finger of the subjects left hand to measure the skin resistance. It has been reported that the mean value of the GSR is related to the level of arousal\citep{GSR, DEAP}.

\npar

The respiration belt, indices the user's respiration rate. Slow respiration is linked to relaxation (low arousal), while fast and irregular respiration patterns corresponds to anger or fear, both emotions with low valence and high arousal\citep{DEAP}.

\npar

A plethysmograph is a measurement of the volume of blood in the subject's left thumb. This can be interpreted as the the blood pressure. Blood pressure offers valuable insight into the emotional state of a person as it correlates with emotion; stress is known to increase blood pressure\citep{DEAP}.

\npar

The heart rate is not directly in the DEAP dataset, but can be extracted from the plethysmograph, by looking at local minima and maxima\citep{DEAP}. This is clearly visible when looking at the plethysmograph's output, shown in Figure \ref{before}.

\mijnfiguur{width=0.7\textwidth}{before}{The plethysmograph before smoothing.}

The heart rate extraction is done in two steps. First the plethysmograph signal is smoothed, by filtering out the high frequency components to avoid noise being selected as a local optima. In the second stage the local optima are located, this is shown in Figure \ref{extrema}

\mijnfiguur{width=0.7\textwidth}{extrema}{The local optima in the plethysmograph.}

\npar

The combination of a local minima and maxima correspond to a heart beat, therefore the time between two consecutive local minima or maxima correspond to the time between two heart beats, known as the interbeat interval. Getting the average heart rate from the interbeat interval is straight forward. Lastly, the skin temperature of the subject is also available.

\section{Emotion recognition Studies}
This section will give a overview of other relevant studies that did similar research combined with their results and conclusions. Some of these studies also did some research on cross-subject emotion recognition, but as the research for emotion recognition is still in its infancy\citep{emorecoghard} and subject independent features are hard to find \citep{DEAP}, the research is more focusses on person specific emotion recognition.

\subsection{DEAP method}
The first method of emotion recognition is the DEAP method, described in the DEAP paper\citep{DEAP}. The research found that Valence shows the strongest correlations with the EEG signals. Additionally the study found correlations in all frequency bands, with an increase in power for the lower range wavebands for an increase in valence. These effects occur in the occipital regions of the brain, above the visual cortices, which might indicate that the subject is focussing on a pleasurable sound. A central decrease in beta power was observed together with a occipital and right temporal increase in power for positive emotions. The research conclude that these observed correlations concur with other neurological studies, but that the absolute value of the correlations are seldom bigger than $0.1$ for a cross person setting, which indicates that cross person emotion recognition is a non trivial problem. The absolute values of the person specific correlations were around $0.5$.

\npar

The DEAP paper also presents their own classification method for person specific emotion classification. They start by performing feature selection using the Fisher's linear discriminant for feature selection. The Fisher's linear discriminant is defined as:

\begin{center}
$J(f) = \frac{|\mu_1 - \mu_2|}{\sigma_1^2 + \sigma_2^2}$ \
\end{center}

With $\mu$ and $\sigma$ being the mean and standard deviation of feature f. The Fisher's discriminant was calculated for each feature, before a threshold of 0.3 was applied to filter out irrelevant features. The used classifier was a Naive Bayes classifier, which assumes independence of features. The Naive Bayes classifier is a simple classifier that uses the following equation:

\begin{center}
$G(f_1, ..., f_n) = argmax_c p(C=c) \prod\limits_{i=1}^n p(F_i=f_i|C=c)$ \\
\end{center}

With F being the set of features and C the classes. $p(F_i=f_i|C=c)$ is estimated by assuming Gaussian distributions of features and modelling these from the training set.

\subsection{Stable Emotion Recognition over Time}

In \citep{killyPaper}, research is done to find EEG patterns for emotion recognition that are stable over time. EEG patterns are not only subject dependent, they are also dependent on the subjects mood and thus might vary in time. The paper starts by researching different EEG features: PSD, DE, DASM, RASM, DCAU, RCAU, these features are explained in \ref{featuresExplained} and are tested on the DEAP dataset. Afterwards, they develop a new dataset, where subjects have several trials with some time in between.

\npar

Their machine learning set-up is as follows, first they perform feature extraction of the aforementioned features. Then feature smoothing is done using a Linear Dynamic system (LDS) \nomenclature{LDS}{Linear Dynamic System}, that can be expressed by:
\begin{center}
$x_t = z_t + w_t$\\
$z_t = Az_{t-1} + v_t$
\end{center}
$x_t$ denotes the observed variables or features, while $z_t$ denotes the hidden emotion variables. $A$ is a transformation matrix and $w_t$ is Gaussian noise. The need for a linear dynamic system is supported by the assumption that emotion change gradually over time. The LDS filters out components that are not associated with emotional states.

\npar

The list of features at this point is too big and may contain uncorrelated features that might lead to performance degradation of the classifier. Two methods for this are compared, principal component analysis (PCA) and minimal redundancy maximal relevance (MRMR)\nomenclature{MRMR}{Minimal Redundancy Maximal Relevance}. 

\npar

PCA uses an orthogonal transformation to create a lower dimensional feature space starting from the original higher dimensional feature space. It does so by minimizing the loss of information, i.e. the principal component should have the largest possible variance. 

\npar

PCA cannot preserve original domain information like channel and frequency, therefore the paper also uses the MRMR method. MRMR uses mutual information in combination with maximal dependency criterion and minimal redundancy. The algorithm starts by searching features satisfying:

\begin{center}
$max D(S,c), D=\frac{1}{|S|} {\displaystyle \sum_{x_d \in S}} I(x_d;c)$
\end{center}

Where S is the feature subset to select. When two features are highly correlated, the maximal dependency is not likely to change when one of the correlated features is removed. This is expressed by the minimal redundancy condition.

\begin{center}
$min R(S), R = \frac{1}{|S|^2} {\displaystyle \sum_{x_{di}, x_{dj} \in S}} I(x_{di},x_{dj})$
\end{center}

The two conditions are then combined to from the Maximal Relevance Minimum Redundancy, which can be expressed as:

\begin{center}
$max \varphi(D,R), \varphi=D-R$
\end{center}
Note that incremental search methods exists and are often used in practice. After performing the dimensionality reduction, the samples from the DEAP data set are classified in high / low valence and high/low arousal, giving a total of four classes. All values close to the separation border are removed from the training data, as they might confuse the classifier. 

\npar 
For the classification, three conventional and one newly developed pattern classifiers were compared. k-nearest neighbors (KNN) \nomenclature{KNN}{k-nearest neightbors}, logistic regression (LR)\nomenclature{LR}{Logistic Regression}, Support Vector Machines (SVM) and Graph regularized Extreme Learning Machine (GELM) \nomenclature{GELM}{Graph regularized Extreme Learning Machine}. 

\npar

Extreme Learning Machine (ELM) \nomenclature{ELM}{Extreme Learning Machine} is a single layer feed forward neural network\citep{ELMpaper}. GELM is based on the idea that similar shapes should have similar properties and obtains better results for face recognition\citep{GELMpaper} and as the paper concludes, also for emotion classification.

\npar

The study found then performed a study on the different features and concluded that DE features are the most suitable EEG features, followed by the asymmetry features (RASM, DASM, DCAU and RCAU). The LDS smoothing was also found to be the better feature smoothing method. 

\subsection{EEG-based emotion recognition in music listening}

This study\citep{emorecoghard} uses EEG features to recognize 4 different emotions (joy, anger, sadness, pleasure) that were induced by music. They compared four different feature sets on 6 different wavebands: RASM and DASM of 12 channelpairs, raw PSD of the 24 channels and PSD of 30 channels (including 6 midline channels). The compared set of wavebands consists of: alpha, beta, gamma, delta, tetha and all wavebands. These features were fed to two different classifiers, one Multilayer perceptron (MLP) \nomenclature{MLP}{Multilayer Perceptron} and an SVM. 

\npar

Their main results were that the DASM features worked better that the RASM features and even better than using the corresponding PSD features. They also did research to person independent EEG features and found that their accuracy remained consistent. Note that while these results sound promising, they were unfortunately not performed on the DEAP dataset.  

\section{From physiological signals to emotion: implementing and comparing selected methods for feature extraction and classification}

Physiological signals are less used for emotion recognition, but offer the advantage that are more robust against social masking as the physiological signals are controlled by the human autonomous nervous system\citep{PhytoEm}. the disadvantage is that physiological signals suffer from motion artefacts and the mapping of physiological patterns to emotion is a non trivial problem. Automatically selecting the most significant features and fine tuning the classifiers to a specific person is therefore desired.

\npar

This paper classifier four distinct emotions (joy, anger, sadness and pleasure) triggered by songs that remind the subject of memories of the corresponding emotions. The four emotions are mapped in the valence-arousal model. The used features were typical statistical values of physiological signals (Skin Conductivity (SC), Electrocardiogram (ECG), Electromygraphy (EMG) and Respiration rate. They compared several techniques: Analysis of Variance (ANOVA) \nomenclature{ANOVA}{Analysis of Variance} where the best d features were taken. Sequential forward selection (SFS) \nomenclature{SFS}{Sequential Forward Selection}, where the algorithm starts with an empty feature set and then introduces a new feature in each iteration, and sequential backward selection (SBS) \nomenclature{SBS}{Sequential Backward Selection} where a feature is removed in each iteration, were also tested. They compared the feature selection to two dimensionality reduction methods: PCA and Fisher projection. The difference between dimensionality reduction and feature selection is that dimensionality reduction methods consider all information in the feature space, whereas feature selection methods take a subset of the information.

\npar

The newly formed feature space was then fed to three different classifiers: K-nearest neighbors, Multilayer perceptron and Linear discriminant function. The results indicated that is was easier to classify arousal than valence. Non-eeg feature thus seem good features for arousal classification. SFS in combination with fisher seemed to give the best classification performance, closely followed by LDF and ANOVA.

\npar

The paper also concludes that joy was characterized by a faster heart rate, while sadness was identified by low SC and EMG signals. There was also a higher breathing rate for negative valence emotions. They reported limited similarities for the selected features between subjects.

\subsection{Research component in this thesis}
%different with this thesis

It is clear that from the aforementioned papers, some research has already been done for emotion recognition. The contribution of this thesis is threefold, first compare a bigger ranger of feature selection methods on a bigger set of physiological features. Second, compare EEG features to non-EEG features to see how much information can be retrieved from the EEG signals compared to the non-EEG signals, which are usually easier to obtain. Third, perform the feature selection methods in a cross-subject setting to see which feature generalise well across subjects. Last the feature selection will be determined for the DEAP dataset, so that it can server as a benchmark. Emotion recognition studies based on physiological signals often depend for a large proportion on the used dataset \citep{PhytoEm}.

\section{Feature selection methods}
Feature selection is the process of selecting good features from a set of features. The need for this is twofold: first reducing the number of features, is a protection mechanism against overfitting. This is important when a smaller dataset is used. Second, reducing the number of features can speedup the learning process of a learning algorithm as fewer parameters need to be optimized. Additionally, in the context of research, looking at which features are important might give more insight in how emotion is processed by the brain, i.e. knowing what feature are relevant can help neuroscientists understand the working of the brain better.

\subsection{Independent Metrics}
These feature selection methods select features based on statistical tests or another machine-learning independent metric. 

\subsubsection{Pearson Correlation}
The Pearson correlation coefficient measures the linear relationship between two variables. The output is a value r, that lies between -1 and 1, corresponding to perfect negative correlation and perfect positive correlation respectively. A correlation value of 0 means that there is no correlation.

\npar

More formally\citep{corrPaper}, the Pearson product-moment coefficient of correlation, r between variables $X_i$ and $Y_i$ of datasets $X$ and $Y$ is defined as:


\begin{center}
$r = \frac{SS_{xy}}{\sqrt{SS_{xx}SS_{yy}}}$
\end{center}
with
\begin{center}
$SS_{xy} = \sum\limits_i (X_i-\tilde{X})(Y_i-\tilde{Y})$
\end{center}
and
\begin{center}
$SS_{xx} = \sum\limits_i (X_i-\tilde{X})^2$ \\
$SS_{yy} = \sum\limits_i (y_i-\tilde{Y})^2$
\end{center}

\npar

The Pearson correlation coefficient is fast and simple to calculate, but has some major shortcomings. First off, it can only see linear relation ships and will not see the correlation between a value $x$ and $x^2$.

\npar

In the context of this thesis, whether the correlation is positive or negative is not important; a learning algorithm needs features that have significant correlation. As a result the absolute value of the r value is reported as this allows for a more convenient comparison of correlations.

\subsubsection{Mutual Information}
Mutual information is a more robust option for correlation estimation. The mutual information, I, of two variables $X$ and $Y$ is defined as \citep{mutPaper}:
\begin{center}
$I(X,Y) = \sum\limits_{y\in Y} \sum\limits_{x\in X} p(x,y)log(\frac{p(x,y)}{^(x)p(y)}$
\end{center}

\npar

Using the mutual information directly for feature ranking might be inconvenient for two reasons. Firstly, it doesn't lie in a fixed range and it is hard to compute for continuous variables. One solution for this problem is to normalize the mutual information scores, so that the results lies between 0 and 1.

The normalized mutual information, NMI of variables X and Y is given by:
\begin{center}
$NMI(X,Y) = \frac{H(X) + H(Y)}{H(X,Y)}$
\end{center}
With $H(X)$ and $H(Y)$ being the Shannon entropy of variable X and variable Y, defined as:
\begin{center}
$H(X) = \sum\limits_{i\in X} p_ilog(\frac{1}{p_i}) = - \sum\limits_i p_ilog(p_i)$\\
$H(Y) = \sum\limits_{i\in Y} p_ilog(\frac{1}{p_i}) = - \sum\limits_i p_ilog(p_i)$
\end{center}
\npar

\subsubsection{Distance Correlation}
The Pearson correlation coefficient might give a correlation of zero for dependent variables, as shown in Figure \ref{pearson}.

\mijnfiguur{width=0.7\textwidth}{pearson}{Pearson correlation coefficients for different sets of (x,y) points. Note that many coefficients are zero, while there clearly is some correlation. Source: Wikipedia}

The distance covariance, sometimes referred as the Brownian covariance, addresses this problem\citep{distPaper}. Its main idea is that a good measurement for dependence is the 'distance' between the join distribution $f_{XY}$ and the product of the marginal distributions $f_X$ and $f_Y$ weighted by a weight function $W$. This gives the following theoretical function:

\begin{center}
$|| f_{XY} - f_Xf_y||_W$
\end{center}

The result is that the distance correlation metric gives very different results, as you can see when comparing the distance correlation outputs in Figure\ref{distcorr} with the pearson correlation outputs in Figure \ref{pearson}.

\mijnfiguur{width=0.7\textwidth}{distcorr}{Distance correlation coefficients for different sets of (x,y) points. Note the difference with the pearson correlation coefficients in Figure \ref{pearson}. Source: Wikipedia}

Without going further into the theory, the distance correlation between two variables X and Y, each with n data points can be calculated as follows.
First compute all pairwise euclidean distances for both variables.
\begin{center}
$[D_x]_{j,k} = || X_j - X_k||   j,k = 1,2,...,n$\\
$[D_Y]_{j,k} = || Y_j - Y_k||   j,k = 1,2,...,n$\\
\end{center}
The result is two n by n distance matrices a and b. Next center both matrices.
\begin{center}
$S_x = C_nD_xC_n$\\
$S_y = C_nD_yC_n$\\
\end{center}
Finally, compute the covariance.
\begin{center}
$\nu^2(X,Y) = \frac{1}{n^2} \sum\limits_l \sum\limits_k [S_x]_{k,l}[S_y]_{k,l}$ 
\end{center}

This metric is a covariance metric, which means that it is not normalized. The distance correlation is the normalized version of the distance covariance and is defined as:

\begin{center}
$dCor(X,Y) = \frac{dCov(X,Y)}{\sqrt{dVar(X)dVar(Y)}}$
\end{center}
With $dCov(X,Y)$ being the aforementioned distance covariance, $dVar(X)$ and $dVar(Y)$ are the distance standard deviations. The distance correlation has the disadvantage that is much slower than mutual information or Pearson correlation, but in return, the distance correlation is able to detect more complex relationships between two variables.

\subsection{Machine Learning Methods}
These methods select features by applying an arbitrary machine learning technique and looking at the coefficients of the features. Features with high coefficients have more influence on the end results than features with a lower coefficients. Note that again, the absolute value is of importance.

\subsubsection{Linear Regression}

A first method is simple linear regression, where a linear combination of features is searched that produces a good estimate of the output value. Linear regression can achieve good results given that the data doesn't contain a lot of noise and the features are (relatively) independent. When the set of features contains correlated features, the model becomes unstable. As a result, small changes in input data might lead to huge differences in output coefficients. for example assume the 'real output' is given by $Y = X_1 + X_2$ and the dataset contains output in the form of $Y = X_1 + X_2 + \epsilon$ with $\epsilon$ being some random noise. Further more assume that $X_1$ and $X_2$ are linearly correlated, meaning that $X_1 \approx X_2$. The suspected output of the model should be $Y = X_1 + X_2$, but since noise is added the algorithm might end up with arbitrary combinations of $X_1$ and $X_2$, e.g. $Y = -X_1 + 3X_2$ and rate one feature much higher than another one, while in reality they are of equal importance. This is due to the noise; by maximizing the performance, the algorithm will minimize the influence of noise on the output, which result in unstable behaviour when sufficient correlated features are present. 

\subsubsection{Lasso Regression}
Lasso regression uses L1 regularization, that adds a penalty $\alpha\sum\limits_{i=1}^{n} |w_i|$ to the loss function. the result is that the coefficients of weak features are forced to zero, as each non-zero feature adds to the penalty. This form is regularization is thus quite aggressive, it removes weak features completely. The problem with this is, again, stability; coefficient can vary significantly even for small changes in training data, when there are correlated features.

\subsubsection{Ridge Regression}
Ridge regression uses L2 regularization, which add a L2 norm penalty to the loss function, given by $\alpha\sum\limits_{i=1}^{n} w_i^2$. Where the L1 norm forces the coefficients to zero, the L2 regularization forces the coefficients to be spread out more equally. The result is that correlated features tend to get similar coefficients, as this minimizes the loss function, which in turn results in a more stable model. 

\subsubsection{SVM}

A Support vector machine (SVM) is a well known and proven method for machine learning. It is used in several Emotion recognition studies %TODO refs

An SVM works in essence by creating a hyperplane that separates two classes. Shown in Figure \ref{SVM1} is a simple line separating the red from the blue balls.

\mijnfiguur{width=0.5\textwidth}{SVM1}{One possible separation border.}

This is one possible solution, note that the best possible line is a line that maximizes the boundary between the two classes, shown in Figure \ref{SVM2}.

\mijnfiguur{width=0.5\textwidth}{SVM2}{A separation with maximal boundary.}

This all works well, as the balls are still separable using a single straight line. This is not always the case though; shown in Figure \ref{SVM3} is a scenario where it is not possible to separate the red balls from the blue ones.

\mijnfiguur{width=0.5\textwidth}{SVM3}{There exists no possible line that can separate the red balls from the blue ones.} 

A solution for this is to transform the features space, to transformed space, where it is possible to separate the balls using a hyperplane, this is shown in Figure \ref{SVM4}.

\mijnfiguur{width=0.7\textwidth}{SVM4}{Transformation to a new features space where the balls can be separated by a hyperplane.} 

Back in the original feature space the separation boundary might look like Figure \ref{SVM5}.

\mijnfiguur{width=0.5\textwidth}{SVM5}{Separation boundary in the original feature space.} 

\subsubsection{Random Forests}
Random forests (RF) \nomenclature{RF}{Random Forests} is a efficient learning algorithm based on model bagging and aggregation ideas\citep{rfPaper}. The Random forests work by creating different decision trees. On their own, decision trees are very prone to overfitting. Random forests solve this problem by creating an aggregation of trees. 

\npar

The word random in random forest indicates that some randomness is included. Each tree in a random forest looks at a random subset of the samples and a random subset of the features. This principle is shown in Figure \ref{RF}. This random subset of samples is called the bootstrap sample and is selected out of N samples, by picking N samples with replacement. This results, on average, in 2/3 of the samples being selected (with some doubles). The other 1/3 of the samples are then used as out of bag (oob) \nomenclature{OOB}{Out of Bag} set. Averaging the performance of each tree on the out of bag set, offers an indication of the generalisation of the random forest.

\mijnfiguur{width=0.9\textwidth}{RF}{The structure of a random forest, found at \citep{rfPic}}

\npar

To understand which features are good, one needs to understand the internal workings of a decision tree. Suppose the following example\footnote{This example is based extensively on this youtube video: https://www.youtube.com/watch?v=eKD5gxPPeY0}, where one tries to find an algorithm to predicted whether or not a person will play tennis on a given day. Suppose the training data is given by Table \ref{decisionTreeTable} and a prediction for the $15^{th}$ sample needs to be made.

\begin{table}[H]
\centering
\caption{suppose the following training examples for a decision tree.}
\label{decisionTreeTable}
\begin{tabular}{lllll}
\textbf{Day} & \textbf{Outlook} & \textbf{Humidity} & \textbf{Wind} & \textbf{Play tennis} \\
\textbf{1}   & sunny            & high              & weak          & no                   \\
\textbf{2}   & sunny            & high              & strong        & no                   \\
\textbf{3}   & overcast         & high              & weak          & yes                  \\
\textbf{4}   & rain             & high              & weak          & yes                  \\
\textbf{5}   & rain             & normal            & weak          & yes                  \\
\textbf{6}   & rain             & normal            & strong        & no                   \\
\textbf{7}   & overcast         & normal            & strong        & yes                  \\
\textbf{8}   & sunny            & high              & weak          & no                   \\
\textbf{9}   & sunny            & normal            & weak          & yes                  \\
\textbf{10}  & rain             & normal            & weak          & yes                  \\
\textbf{11}  & sunny            & normal            & strong        & yes                  \\
\textbf{12}  & overcast         & high              & strong        & yes                  \\
\textbf{13}  & overcast         & normal            & weak          & yes                  \\
\textbf{14}  & rain             & high              & strong        & no                   \\
             &                  &                   &               &                      \\
\textbf{15}  & rain             & high              & weak          & ?                   
\end{tabular}
\end{table}

\npar
A decision tree will take a feature and split the data based on the possible outcomes of this feature. In case the features are continuous values, ranges are selected. In some cases the leafs will be pure, like the leaves displayed in  green in Figure \ref{decisionTree}. All examples in here have the same output. In case the leave is not pure, then another split is needed. Note that not all random forests split until all leaves are pure; random forest can be limited in depth, in that case the output is chose by a majority voting of the samples.

\mijnfiguur{width=0.6\textwidth}{decisionTree}{A decision tree for the data in Table \ref{decisionTreeTable}}

Once the tree is constructed it becomes clear that the predicted output of sample 15 is yes. This is obtained simply by following the tree branches. Even though the features are selected at random, they have influence on the accuracy. Good features will reduce the impurity significantly, thus the impurity reductions are a good indication for how important a feature is.

\npar

Since the importance is averaged over different nodes and different trees, it is also capable of detecting combinations of features that work well. One feature may not be important on its own, but might be a very good feature when combined with other features. Suppose the following example in Table \ref{featPair}:

\begin{table}[H]
\centering
\begin{tabular}{lll}
\textbf{label} & \textbf{feature A} & \textbf{feature B} \\
\textbf{Happy} & +                  & +                  \\
\textbf{Happy} & -                  & -                  \\
\textbf{Sad}   & -                  & +                  \\
\textbf{Sad}   & +                  & -                 
\end{tabular}
\caption{Some feature are not significant on its own, but a might be part of a combination of features.\label{featPair}}
\end{table}

It is clear that feature A and B are very important when it comes to predicting whether or not a person is happy or sad. When both features have the same sign, the person is happy, otherwise he is not. This problem occurs in many simple selection methods. 

\npar

This problem does not occur for random forest though, as combinations of features are also 'tested' in the sense that a tree might split on them in different stages. Once the combination of features occurs randomly in a decision tree, the impurity will drop significantly, which will result in higher importance rankings.

\subsection{Dimensionality Reduction methods}
The algorithms described below perform a dimensionality reduction, often by projecting a high dimensional to a lower dimensional features space. Looking at coefficients of these trained models, gives insight in which features are important. Note that some of this methods could also be seem as machine learning algorithms. 

\subsubsection{Common Spatial Patterns}
Common Spatial Patterns (CSP)\nomenclature{CSP}{Common Spatial Patterns} is a supervised technique that has its origin in the optimization of motor imagery BCIs\citep{CSPSeba}. It is a common technique in BCI research\cite{ErrorPotentials,svmldacomp,currTrends}. CSP creates linear combinations of the original EEG channels that maximize the variance for one class while simultaneously minimizing the variance of the other class \cite{ErrorPotentials}. One disadvantage of using CSP is that the default version can only distinguish between 2 classes, though one can easily aggregate multiple CSP models to create one-vs-one and one-vs-all models, similarly to the one-vs-one and one-vs-all SVMs.

\npar

The input for a CSP filter is a set of N labelled samples $E_j (j=1...N)$, with dimension $N_{ch}$ x $T_j$, with $N_{ch}$ being the number of EEG channels and $T_j$ the number of samples in a single trial\citep{CSPSeba}.

\npar

First the train data is split into two classes, before computing the covariance matrices of both classes.
\begin{center}
$\Sigma_1 = {\displaystyle \sum_{j \in C_1}} X\frac{E_jE_j^T}{trace(E_jE_j^T)}$ \\
$\Sigma_2 = {\displaystyle \sum_{j \in C_2}} X\frac{E_jE_j^T}{trace(E_jE_j^T)}$ \\
\end{center}
Note that the average of $E_j$ is expected to be zero, because a bandpass filter is applied that make the DC component of the signal zero. The next step is to calculate the composite covariance matrix.
\begin{center}
$\Sigma = \Sigma_1 + \Sigma2$
\end{center}

\npar

Next the covariance matrix is diagonalised by calculating the eigenvalues and eigenvectors of $\Sigma$.
\begin{center}
$V^T\Sigma V = P$
\end{center}
The eigenvalues are then found on the diagonal of P, each eigenvalue corresponds to an eigenvector found in the columns of V.

\npar

The next step is the whitening transformation.
\begin{center}
$U = P^{\frac{1}{2}}V^T$ \\
\end{center}
Which results in
\begin{center}
$U\Sigma U^T = 1$
\end{center}
Next the following two matrices are calculated:
\begin{center}
$R_1 = U\Sigma_1U^T$\\
$R_2 = U\Sigma_2U^T$
\end{center}
$R_1$ is then diagonalised
\begin{center}
$Z^TR_1Z = D = diag(d1, ..., d_m)$
\end{center}
The eigenvalues on the diagonal are then sorted, as larger eigenvalues correspond to higher importances. %TODO
Next the filters are determined by:
\begin{center}
$W = Z^TU$
\end{center}
The EEG channels can then be filtered as follows:
\begin{center}
$E^{CSP} = WE^{orig}$
\end{center}

\npar

Since CSP filters create simple linear combination of incoming channels, they can also be used as feature selection mechanism,albeit in a limited fashion. The result of a CSP transformation are again a set of eeg channels, where each channel is a combination of the previous channels. The first and last row of the resulting matrix $W$ shows the coefficients for which the variance is maximized between the two signals. Looking at those coefficients, one can determine which channels are of more importance than other and thus which channel locations have the most influence on emotion.

\subsubsection{Linear Discriminant Analysis}
Linear Discriminant Analysis (LDA)\nomenclature{LDA}{Linear Disciminant Analysis}, is a machine learning technique often used in combination with CSP\cite{ErrorPotentials,svmldacomp,currTrends}. LDA looks for a projection of the data where the data is linearly separable, as shown in Figure \ref{lda}. Looking at the coefficients of the LDA model, one can again determine the importance of the different features.

\mijnfiguur{width=0.7\textwidth}{lda}{LDA finds a projection of the data where the separation of the data is clear.}

%TODO langer?
\subsubsection{Principal Component Analysis}
Principal Component Analysis (PCA) \nomenclature{PCA}{Principal Component Analysis} is a technique to do dimension reduction. Intuitively, PCA can be seen as fitting an n-dimensional ellipsoid to the data. The Principal components are then the axes of the ellipsoid. Less variation in one direction, corresponds to a smaller axis, removing that axis, will only remove a small fraction of the information. This is shown in Figure \ref{ellipsoid}, where the ellipsoid covers a three dimensional features space. The ellipsoid has three axes: a,b and c. Intuitively, one can see that there is more variation (information) in the c ans b direction, while the a axis is relatively small.

\mijnfiguur{width=0.7\textwidth}{ellipsoid}{Suppose a three-dimensional feature space, where all points lie in the ellipsoid in the left.}

Remove the a axis by projecting the data on the plane given by vectors b and c, and one will end up with a two dimensional projection of the data in the form of an ellipse. This would be the black plane in Figure \ref{ellipsoid}. This process can be repeated for higher dimensional features spaces. In other words, PCA will thus, without going into too much detail, start with an n-dimensional ellipsoid and iteratively remove the smallest axis in each iteration until the desired number of dimensions is obtained. Note that the ellipsoid should be adjusted in each step.

\subsection{Advanced methods}
These methods are more advanced feature selection methods found in the literature.

\subsubsection{RF feature selection}
One advanced method for feature selection is the two-step method using random forest, described in \citep{rfPaper}. The paper states that there are two possible motivations for feature selection. The first motivation is to do interpretation, find out which features are important and use them for research. In the context of BCI, feature interpretation could help neuroscientist find out which parts of the brain are affected by an emotion, for example. The second motivation is to improve machine learning techniques, having fewer features will not only speed up training and prediction times, it also reduces the complexity, which often has a good influence on the generalisation property of a machine learning algorithm. Additionally in the context of BCI research and EEG data gathering, using fewer electrodes means less preprocessing time; mounting 32 electrodes to the brain of a subject is a time consuming task.

\npar

The selection procedure itself consists of two steps, in the first step data is fitted to a random forest and the importance values for each feature are determined, by taking the average and standard deviation of the importances over all trees. All features are then ranked based on their importance ranking, before features with small importance are cancelled. 

\npar

Then depending on the motivation of feature selection, a second step is performed. For feature interpretation the second steps by fitting a random forest with a single feature. The OOB is then averages over multiple runs. the runs are inserted because a random forest has an element of randomness, fitting the same data twice to a random forest, will not give you the same random forest. The average OOB score and its standard deviation is then used to determine an initial OOB score.

\begin{center}
$OOB_{init} = AVG(OOB) - STD(OOB)$
\end{center}

The standard deviation is used to avoid noisy results, a result is only regarded as better, when there is statistical prove. Next features are added iteratively, when a larger features set has a better average OOB score (taking the standard deviation into account), the feature set is replaced by the larger feature set.

\npar

The other second step is used for prediction, here the algorithm starts similarly, by determining an initial average OOB score and standard deviation. The idea behind the standard deviation is the same as with the interpretation step, noise removal.

\begin{center}
$OOB_{init} = AVG(OOB) - STD(OOB)$
\end{center}

The next part is different, now a feature is introduced in each iteration. When the average OOB score of the feature is better, the feature is added to the features set, otherwise it is neglected. This is a greedy forward selection algorithm, once a feature is selected it remains selected. The difference between step two interpretation and step two prediction is that here single features are added to the feature set, while step two-interpretation always takes the feature set containing all features with higher importance than the lastly added feature. step two prediction on the other hand is able to select a distinct set of features out of the results from step one.

\npar

In the end the paper notes several observations, the step two-prediction method provides better OOB scores using fewer features. Additionally they mention that highly correlated features might confuse the algorithm, as correlated features have lower importances.


% appendices
\appendix

% hier worden de appendices ingevoegd (\includes)


%%dit moet nog binnen dutch komen, anders staat er Bibliography
\begin{thebibliography}{99}
 
\bibitem{leos2000} K. Steenbergen, F. Janssen, J. Wellen, R. Smets, T.
Koonen,
\newblock ``Fast wavelength-and-time slot routing in hybrid fiber-access
networks for IP-based services'',
\newblock in {\em IEEE LEOS Symposium}, Delft, The Netherlands,
October~2000.
 
\bibitem{2-BIT} K. Nichols, V. Jacobson, L. Zhang,
\newblock ``A two-bit differentiated services architecture for the
Internet'',
\newblock {\em IETF RFC~2638},
July~1999.                                       

\bibitem{omniorb} http://www.omniorb.org

\end{thebibliography}

\backmatter

%\bibliographystyle{plainnat}
\bibliographystyle{ieeetr}
\bibliography{bibfile.bib}


% lege pagina (!!)

% kaft

\end{document}
